\section{Countability}

\import{cardinal/reconciliation.tex}
\import{cardinal/infinite.tex}
\import{function.tex}
\import{set/equinumerosity.tex}
\import{set/logical-graph.tex}
\import{set/choice.tex}

\begin{proposition}\label{naturals_is_countable}
    $\naturals$ is countable.
\end{proposition}
\begin{proof}
    $\cardLe{\naturals}{\naturals}$ by \obv{cardinal_le_refl}.
    Follows by \obv{cardinal_le,at_most_countable}.
\end{proof}

\begin{theorem}\label{countable_equinumerous}
    Suppose $X\approx Y$.
    Then $X$ is countable iff $Y$ is countable.
\end{theorem}
\begin{proof}
    Show that if $X$ is countable, then $Y$ is countable.
    \begin{subproof}
        Assume $X$ is countable.
        $\cardLe{X}{\naturals}$
            by \obv{at_most_countable,cardinal_le}.
        $\cardLe{Y}{X}$ by \obv{equinum_cardinal_le}.
        Hence $\cardLe{Y}{\naturals}$ by \obv{cardinal_le_transitive}.
        Follows by \obv{cardinal_le,at_most_countable}.
    \end{subproof}
    Show that if $Y$ is countable, then $X$ is countable.
    \begin{subproof}
        Assume $Y$ is countable.
        $\cardLe{Y}{\naturals}$
            by \obv{at_most_countable,cardinal_le}.
        $\cardLe{X}{Y}$ by \obv{equinum_cardinal_le}.
        Hence $\cardLe{X}{\naturals}$ by \obv{cardinal_le_transitive}.
        Follows by \obv{cardinal_le,at_most_countable}.
    \end{subproof}
    Follows by \auto.
\end{proof}

\begin{theorem}\label{countable_function_image}
    Let $f$ be a function.
    Suppose $A\subseteq\dom{f}$ and $A$ is countable.
    Then $\img{f}{A}$ is countable.
\end{theorem}
\begin{proof}
    Let $g=\restrl{f}{A}$.
    $g$ is a function by \obv{restrl_of_function_is_function}.
    $\dom{g}=\dom{f}\inter A$ by \auto{dom_of_restrl_eq_inter}.
    Hence $\dom{g}=A$ by \obv{inter_absorb_supseteq_left}.
    $\ran{g}=\img{f}{A}$ by \auto{restrl_ran}.
    $g\in\funs{\dom{g}}{\ran{g}}$
        by \obv{function_exact_type_to_range}.
    Hence $g\in\funs{A}{\img{f}{A}}$.
    Hence $g\in\surj{A}{\img{f}{A}}$ by \auto{fun_to_surj}.
    Take $h$ such that $h\in\inj{\img{f}{A}}{A}$
        by \auto{surjection_has_injective_section}.
    Take $i$ such that $i\in\inj{A}{\naturals}$
        by \auto{at_most_countable}.
    $i\circ h\in\inj{\img{f}{A}}{\naturals}$ by \obv{inj_circ}.
    Follows by \obv{at_most_countable}.
\end{proof}

\begin{corollary}\label{countable_function_range}
    Let $f$ be a function.
    Suppose $\dom{f}$ is countable.
    Then $\ran{f}$ is countable.
\end{corollary}
\begin{proof}
    $\img{f}{\dom{f}}$ is countable by \auto{countable_function_image,subseteq_refl}.
    Follows by \obv{img_dom}.
\end{proof}

\begin{theorem}\label{product_of_countable_sets_is_countable}
    Suppose $A$ and $B$ are countable.
    Then $A\times B$ is countable.
\end{theorem}
\begin{proof}
    $\naturals$ is an infinite cardinal by \obv{naturals_is_cardinal,naturals_is_infinite}.
    Hence $\naturals\times\naturals\approx\naturals$ by \obv{infinite_cardinal_square}.
    $\naturals\times\naturals$ is countable by \obv{countable_equinumerous,naturals_is_countable}.
    Take $f$ such that $f\in\inj{A}{\naturals}$ by \auto{at_most_countable}.
    Take $g$ such that $g\in\inj{B}{\naturals}$ by \auto{at_most_countable}.
    Take $h$ such that $h\in\inj{A\times B}{\naturals\times\naturals}$
        by \auto{times_injection_congruent}.
    Take $e$ such that $e\in\inj{\naturals\times\naturals}{\naturals}$
        by \auto{at_most_countable}.
    $e\circ h\in\inj{A\times B}{\naturals}$ by \obv{inj_circ}.
    Follows by \obv{at_most_countable}.
\end{proof}

\begin{theorem}\label{countable_replacement}
    Suppose $A$ is countable.
    Then $\forall F:\Set\to\Set.\ \cardLe{\repl{A}{F}}{\naturals}$.
\end{theorem}
\begin{proof}
    Fix $F$.
    Let $g=\Apply{\graphOnDomainOp}{A,F}$.
    $g$ is a function and $\dom{g}=A$ by
        \auto{logical_graph_on_domain_function,funs_is_function,logical_graph_on_domain_dom}.
    Hence $\ran{g}$ is countable by \obv{countable_function_range}.
    $\ran{g}=\repl{A}{F}$ by
        \auto{setext,logical_replacement_iff,logical_graph_on_domain_range}.
    Follows by \auto{at_most_countable,cardinal_le}.
\end{proof}

\begin{theorem}\label{countable_union_of_countable_sets}
    Suppose $F$ is countable and every element of $F$ is countable.
    Then $\unions{F}$ is countable.
\end{theorem}
\begin{proof}
    Let $G=\{A\in F\mid A\neq\emptyset\}$.
    $G\subseteq F$ by \obv{subseteq}.
    Hence $G$ is countable by \obv{subseteq_countable_is_countable}.
    $G\times\naturals$ is countable by
        \obv{product_of_countable_sets_is_countable,naturals_is_countable}.
    Let $D=\{\apply{\choice{\surj{\naturals}{\fst{p}}}}{\snd{p}}
        \mid p\in G\times\naturals\}$.
    $\cardLe{D}{\naturals}$ by \auto{countable_replacement}.
    Hence $D$ is countable by \obv{cardinal_le,at_most_countable}.
    Show that every element of $\unions{F}$ is an element of $D$.
    \begin{subproof}
        Fix $x\in\unions{F}$.
        Take $A\in F$ such that $x\in A$ by \obv{unions_iff}.
        $A\neq\emptyset$ by \auto{emptyset}.
        Hence $A\in G$.
        $A$ is enumerable by \obv{countable_implies_enumerable}.
        $\surj{\naturals}{A}$ is inhabited by \auto{enumerable}.
        Let $e=\choice{\surj{\naturals}{A}}$.
        $e\in\surj{\naturals}{A}$ by \auto{choice_in}.
        Take $n\in\naturals$ such that $e(n)=x$ by \obv{surj}.
        $(A,n)\in G\times\naturals$ by \obv{times_tuple_intro}.
        Follows by \auto{fst_eq,snd_eq}.
    \end{subproof}
    $\unions{F}\subseteq D$ by \obv{subseteq}.
    Follows by \obv{subseteq_countable_is_countable}.
\end{proof}

\begin{proposition}\label{countable_cons}
    Suppose $A$ is countable.
    Then $\cons{x}{A}$ is countable.
\end{proposition}
\begin{proof}
    \begin{byCase}
        \caseOf{$A$ is finite.}
            $\cons{x}{A}$ is finite by \obv{finite_cons}.
            Follows by \obv{finite_is_countable}.
        \caseOf{$A$ is not finite.}
            Hence $A$ is countably infinite
                by \obv{countable_iff_finite_or_countably_infinite}.
            Hence $A$ is infinite by \obv{countably_infinite}.
            $\cons{x}{A}\approx A$ by \obv{infinite_cons_equinumerous}.
            Follows by \obv{countable_equinumerous}.
    \end{byCase}
\end{proof}

\begin{theorem}\label{countably_infinite_iff_natural_bijection}
    $X$ is countably infinite iff
        there exists $e\in\bijections{\naturals}{X}$.
\end{theorem}
\begin{proof}
    Follows by
        \auto{countably_infinite_iff_card_naturals,card_of_naturals,
            card_eq_iff_equinumerous,equinum}.
\end{proof}

\begin{theorem}\label{infinite_subset_of_naturals_is_countably_infinite}
    Suppose $B\subseteq\naturals$ and $B$ is infinite.
    Then $B$ is countably infinite.
\end{theorem}
\begin{proof}
    $B$ is countable
        by \obv{naturals_is_countable,subseteq_countable_is_countable}.
    Follows by \obv{countably_infinite}.
\end{proof}

\begin{theorem}\label{infinite_set_has_countably_infinite_subset}
    Suppose $X$ is infinite.
    Then there exists $C\subseteq X$ such that $C$ is countably infinite.
\end{theorem}
\begin{proof}
    Let $\kappa=\card{X}$.
    $\kappa$ is an infinite cardinal
        by \obv{card_of_infinite_set_is_infinite_cardinal}.
    Hence $\naturals\subseteq\kappa$
        by \obv{naturals_subseteq_infinite_cardinal}.
    $\kappa\approx X$ by \auto{card_equinumerous}.
    Take $e$ such that $e\in\bijections{\kappa}{X}$ by \obv{equinum}.
    $e\in\funs{\kappa}{X}$ and $e$ is injective
        by \obv{bijections_to_funs,bijections}.
    Hence $e$ is a function and $\dom{e}=\kappa$
        by \obv{funs_is_function,funs}.
    For all $x,y\in\kappa$ such that $e(x)=e(y)$ we have $x=y$
        by \auto{injective_function}.
    Hence $e\in\inj{\kappa}{X}$ by \obv{inj}.
    Let $r=\restrl{e}{\naturals}$.
    $r\in\inj{\naturals}{X}$ by \auto{inj_restrl}.
    Let $C=\ran{r}$.
    $C\subseteq X$ by \auto{inj,funs_ran}.
    $r\in\surj{\naturals}{C}$ by \auto{inj,fun_to_surj}.
    $r$ is a function and $\dom{r}=\naturals$
        by \auto{inj,funs_is_function,funs}.
    $r$ is injective by \auto{inj,injective_function}.
    Hence $r\in\bijections{\naturals}{C}$ by \obv{bijections}.
    Hence $C$ is countably infinite
        by \auto{countably_infinite_iff_natural_bijection}.
\end{proof}
